\chapter{Quantenobjekte}


\margintext{{\em 2024-10-30}}
\section{Quantenobjekte am Doppelspalt}

Erfolgreich (Interferenzmuster) durchgeführt:

\startitemize[packed]
\item Elektronen
\item Neutronen
\item Fullerene
\stopitemize


\startframedtask
Wie erzeugen einzelne Moleküle ein Interferenzmuster?
\stopframedtask

Das interferenzmuster entsteht ebenfalls wenn jeweils {\bf einzelne} Elektronen
durch den Doppelspalt geschossen werden. Dafür muss aber eine taugliche Anzahl
an Teilchen druch den Doppelspalt geschickt werden. Sollten es zu wenig sein,
ist das Interferenzmuster nicht deutlich und die verteilung sieht zufällig aus.

\startitemize[packed]
\item einzelnes Elektron landet an zufälligem Ort
\item mit ausreichend vielen Elektronen erhält man ein erkennbares Interferenzmuster 
\item funktioniert auch für andere Quantenobjekte
\stopitemize

\startdefinition[title={Quantenobjekte}]

Quantenobjekte sind alle Objekte, die \quotation{klein genug} sind, um
prinzipiell beim Durchgang durch einen Doppelspalt / ein Gitter ein
Interferenzmuster zu zeigen.

\stopdefinition

Man kann hier also stochastische Vorhersagen über die Auftreffwahrscheinlichkeit der Elektronen an den uterschiedlichen Stellen des Schirmes machen.

Jedes Elektron trifft an exakt einem Ort auf.

Wenn man Messen will, durch welchen Spalt die Elektronen sich bewegt haben, ist
kein Interferenzmuster mehr sichtbar.

Auch wenn man einen der Spalt misst, nachdem das Elektron schon längst durch
sein müsste, \quotation{weiß} das Elektron das ebenfalls.

Liste Eigenschaften:

\startitemize[packed,m]
\item Stochastische Deutung:\crlf Einzelergebnisse sind zufällig, man kann nur Wahrscheinlichkeiten angeben.
\item Interferenz:\crlf Quantenobjekte können Interferenz zeigen.
\item Determiniertheit:\crlf Einzelergebnisse sind nach der Messung eindeutig.
\item Komplementarität:\crlf Messungen können andere Messungen verhindern bzw. weniger genau machen. Messungen des durchquerten Spaltes im Doppelspalt verändert Interferenz
\stopitemize
